% Az OTDK-ra (38. OTDK Informatikatudományi Szekció) beküldendő kötelező függelékek.
% A kari TDK-változatból kikerültek; az OTDK-változatban a fuggelek.tex után kell beilleszteni.

% =====================================================================
%  A 38. OTDK Informatikatudományi Szekció kötelező elemei (1. számú melléklet)
% =====================================================================

\chapter{Rezümé}
\label{app-rezume}

\noindent{\small\itshape A 38.~OTDK Informatikatudományi Szekció 1.~számú melléklete szerinti rezümé.}

\medskip
A feleletválasztós vizsgakérdések írása az oktatói munka egyik legidőigényesebb része. A felhőben futó nyelvi
modellek gyors segítséget adnak, de nem ellenőrzik, hogy a helyes választ a tananyag alátámasztja-e, nem
szabályozzák a vizsga szerkezetét, és a tananyagot harmadik félhez juttatják. Helyben futó kis modellekkel az
adatvédelmi probléma megszűnik, de nem ismert, melyik generálási megoldás ad velük megbízható vizsgát.
Kutatási kérdésünk, hogy a feladat ellenőrizhető lépésekre bontása lehetővé teszi-e ezt, és mitől függ az
eredmény.

Lépésekre bontott generálási módszert javasoltunk (\Mimir{}), és kutatási prototípusként egy saját, nyílt
kiértékelő keretrendszerben valósítottuk meg. A keretrendszer 14 konfigurációja egyszerre egy tényezőben tér
el: vizsgaterv, kérdésenkénti keresés és forráshivatkozás, ellenőrző lánc, fogalomgráf, hibrid keresés és
darabolás. A kérdéseket forrásalapú mutatókkal és egy független, bíróként pontozó nyelvi modellel értékeltük,
a konfigurációkat dokumentumszintű páros próbákkal hasonlítottuk össze. Három vizsgálatban, legfeljebb 30
magyar és angol dokumentumon, három fogyasztói GPU-n futó, 4 bites modellel mértünk.

Saját eredményünk a \Mimir{} generálási módszere, a kiértékelő keretrendszer, a vizsgálatok elvégzése és
elemzése, valamint a belőlük levont következtetések. Kimutattuk, hogy a naiv visszakeresés konkrét okból
törékeny: a feladatleírást használja lekérdezésként, és rövid, tartalom nélküli darabokat talál meg. A
lépésekre bontás gyenge modellnél nélkülözhetetlen ($0{,}28\to0{,}67$), a Qwen2.5-7B vele a szervermodell
teljes dokumentumos promptjának minőségét érte el, az erős Gemma 3 12B-nél pedig a rögzített darabolás
önmagában is elég volt. \tbd{[A témavezető hozzájárulása egy mondatban.]}

Fogyasztói GPU-n a vizsgakérdések minőségét elsősorban a választott modell és a visszakeresés lekérdezése
határozza meg; a lépésekre bontás gyenge modellnél a minőséget, erős modellnél az ellenőrizhetőséget adja. A
mérések alapján ezért a vizsgaterv kérdésenkénti forráshivatkozással, szabályalapú ellenőrzéssel és rögzített
darabolással ajánlott. Módszertani tanulság, hogy az ellenőrző és a bíró ítélete véletlenszerűen egyezett, és
több csendes hiba, köztük a válaszszivárgás, a valósnál jobbnak mutatta volna az eredményt.

\medskip
\noindent\textbf{Kulcsszavak:} feleletválasztós teszt, kvantálás, adatvédelem, oktatástechnológia, magyar
nyelvtechnológia

% ---------------------------------------------------------------------
\chapter{Nyilatkozat a mesterséges intelligencia alapú eszközök használatáról}
\label{app-mi-nyilatkozat}
\chaptermark{MI-használati nyilatkozat}
% ---------------------------------------------------------------------

A 38.~OTDK Informatikatudományi Szekció 1.~számú melléklete szerint nyilatkozunk a mesterséges intelligencia
(MI) alapú eszközök használatáról, a tároló verziótörténete alapján; a pirossal jelölt részeket a szerzők
pontosítják. A vizsgált nyelvi modellek és a bíró a kutatás tárgyai, illetve mérőeszközei
(\ref{ch:kiertekeles}.~fejezet), nem a dolgozat elkészítésének segédeszközei; az általuk írt kérdések csak
idézett példaként szerepelnek a szövegben.

\begin{sloppypar}
\begin{description}[style=nextline,leftmargin=0.6cm]
\item[Eszköz] A Claude (Anthropic) kódolási asszisztens (Claude Code); a commitok \code{Co-Authored-By} sora
szerint a Claude Opus 5.5 modell. \tbd{[Ha más modellváltozat vagy felület (például webes csevegőfelület) is
szerepet kapott, itt kell feltüntetni.]}
\item[Szoftverfejlesztés] A kiértékelő keretrendszer (\code{tests/eval}, \code{tests/new~eval}) és az elemző
szkriptek (\code{tdk/elemzes}) fejlesztése: kódírás, hibakeresés, egységtesztek; a tároló dolgozaton kívüli
részeinél is. \tbd{[A szerzők pontosítják, mely részeket írták saját kezűleg, és hogyan ellenőrizték a
generált kódot (kódátnézés, tesztek, kézi futtatás).]}
\item[A dolgozat szövege] A fejezetek megfogalmazása és átdolgozása, a LaTeX-forrás, a diagramforrások és az
ábrákat előállító szkriptek. \tbd{[A szerzők pontosítják, mely fejezeteket és szakaszokat írták saját
kezűleg, melyeket fogalmazott meg az asszisztens az ő irányításukkal, és milyen ellenőrzésen ment át a
szöveg.]}
\item[Mérték] \tbd{[A használat mértéke fejezetenként vagy összesítve, a szerzők becslése szerint.]}
\item[Más eszközök] \tbd{[Más MI-alapú eszközök, ha voltak (például fordító, helyesírás-ellenőrző,
képgenerálás, más csevegő asszisztens), és felhasználásuk módja.]}
\item[Ellenőrzés] A dolgozat számai a nyers kimenetekből egyetlen szkripttel (\code{tdk/elemzes/elemzes.py})
készültek, nem nyelvi modell által írt szövegből; csak az irodalomjegyzékben felvett forrásokra hivatkozunk.
\tbd{[A szerzők megerősítik, hogy a hivatkozott forrásokat és az MI-vel írt részeket ellenőrizték.]} A
tartalomért és a következtetésekért a szerzők felelnek. \tbd{[Hely, dátum, a szerzők aláírása.]}
\end{description}
\end{sloppypar}

% ---------------------------------------------------------------------
\chapter{Eredménynyilatkozat}
\label{app-eredmeny}
% ---------------------------------------------------------------------

A 38.~OTDK Informatikatudományi Szekció 1.~számú melléklete szerint az alábbiakban elhatároljuk a szerzők
saját eredményeit a témavezető és a tágabb kutatócsoport eredményeitől. A szerzők közötti munkamegosztást és a
témavezető szerepét a pirossal jelölt helyeken a szerzők töltik ki.

\section{A szerzők saját eredményei}

\paragraph{A generálási módszer.} A \Mimir{} módszer megtervezése és kutatási prototípusként való
megvalósítása: a darabolás, a vizsgaterv és a Bloom-szintek kiosztása, a kérdéshelyenkénti keresés, a
forrásrészletekre hivatkozó egyenkénti generálás, az ellenőrző lánc, a szivárgás elleni szabályok és a
fogalomgráf, valamint a helyi futtatás GPU-memóriaigényének modellje és mérése (\ref{ch:modszer}.~fejezet).
\tbd{[Munkamegosztás a szerzők között.]}

\paragraph{A kiértékelő keretrendszer.} Az egytényezős konfigurációk, a forrásalapú mutatók (szivárgás,
lefedettség, vak megoldhatóság), a bíró bekötése, az oktatói értékelési protokoll és a dokumentumszintű
statisztikai elemzés, valamint a 30 dokumentumos adatkészlet összeállítása a két nyílt adatkészletből és a saját
szövegekből (\ref{ch:kiertekeles}.~fejezet). \tbd{[Munkamegosztás a szerzők között.]}

\paragraph{A vizsgálatok és az elemzés.} A pilot, az előzetes és a fő vizsgálat megtervezése és futtatása, a
nyers kimenetek kézi átolvasása, a mérési hibák (válaszszivárgás, csendes tartalékágak, hiányzó bírói
ítéletek) feltárása, valamint az eredmények elemzése és a dolgozat összes számát a nyers kimenetekből
előállító elemző szkript (\ref{ch:eredmenyek}. és \ref{ch:modellek}.~fejezet). \tbd{[Munkamegosztás a szerzők
között.]}

\paragraph{Az értékelés.} A hipotézisek értékelése, az eredmények összevetése az irodalommal, az érvényesség
korlátainak elemzése és a mérések alapján ajánlott konfiguráció GPU-szintenként (\ref{ch:diszkusszio}.~fejezet).

\section{A témavezető és a tágabb kutatócsoport hozzájárulása}

\tbd{[A témavezető szerepe (például a téma kijelölése, konzultációk, a dolgozat átnézése, az egyetemi
GenAI-szerver hozzáférésének biztosítása).]} \tbd{[A tágabb kutatócsoport vagy más hallgatók hozzájárulása,
ha volt; ha nem volt, ennek kimondása.]} \tbd{[Az oktatói értékelésben részt vevő oktatók szerepe
(\ref{sec:oktatoi}.~szakasz).]}

\section{Átvett, nem saját eredmények}

Nem saját eredmény a vizsgált nyelvi modellek (Qwen2.5-7B, PULI-LlumiX 8B, Gemma 3 12B, Llama-3.1-8B,
gpt-oss-120b, a szervermodell) és a beágyazó modell (multilingual-e5-base); az EduQG és a MILQA adatkészlet,
amelyből a 27 nyílt dokumentum és a referenciakérdések származnak; a felhasznált nyílt forrású szoftverek
(például az Ollama, a \texttt{llama.cpp} és a Qdrant); valamint az egyetemi GenAI-szerver, amelyen a
szervermodell és a bíró futott. A kódolási asszisztens használatát \az{\ref{app-mi-nyilatkozat}}.~függelék írja
le.

% ---------------------------------------------------------------------
\chapter{Az adatok forrása és nyilvánossága}
\label{app-adatok}
% ---------------------------------------------------------------------

Az adatkészlet összetételét \az{\ref{sec:adatkeszlet}}.~szakasz írja le; a dokumentumok licencét az
adatkészlet leírófájljai (\code{dataset/manifest*.yaml}) is rögzítik.

\begin{table}[htbp]
\centering
\caption[Az adatok és a modellek forrása, licence és nyilvánossága.]{Az adatok és a modellek forrása, licence
és nyilvánossága. A dokumentumok a fő vizsgálat 30 dokumentuma; az előzetes vizsgálat hat dokumentuma ezek
közül való. A helyi modellek 4 bites (Q4\_K\_M) kvantálással futottak az Ollamán; a modellek licencét a
modellkártyák alapján még ellenőrizni kell.}
\label{tab:adatforrasok}
\small
\begin{tabularx}{\textwidth}{@{}>{\raggedright\arraybackslash}p{3.3cm}>{\raggedright\arraybackslash}X>{\raggedright\arraybackslash}p{3.6cm}@{}}
\toprule
Adat vagy modell & Forrás & Licenc, nyilvánosság \\
\midrule
3 saját magyar szöveg & a szerzők saját szövegei (\code{tests/TestDocs}) & saját; a tárolóban nyilvános \\
3 angol tankönyvi szöveg & OpenStax-tankönyvek az EduQG adatkészletből~\cite{hadifar2023eduqg} & CC BY 4.0 \\
24 magyar szöveg & magyar Wikipédia-cikkek a MILQA adatkészletből~\cite{novak2023milqa} & CC BY-SA 4.0 \\
referenciakérdések & EduQG és MILQA; a saját szövegekhez saját & a forrás licence; a tárolóban \\
nyers eredmények & saját futások (\ref{app-reprodukcio}.~függelék) & \tbd{[MIT-licenc alá tartoznak-e]}; a tárolóban \\
oktatói értékelések & részt vevő oktatók & \tbd{[nyilvánosság; az értékelők adatai nélkül]} \\
\midrule
Qwen2.5-7B & generátor (pilot, előzetes)~\cite{qwen25}; \code{qwen2.5:7b} & \tbd{[várhatóan Apache 2.0]} \\
PULI-LlumiX 8B & generátor (fő)~\cite{pulillumix}; \nolinkurl{hf.co/mradermacher/PULI-LlumiX-Llama-3.1-GGUF} & \tbd{[várhatóan Llama 3.1 Community License]} \\
Gemma 3 12B & generátor (fő)~\cite{gemma3}; \code{gemma3:12b} & \tbd{[várhatóan Gemma Terms of Use]} \\
Llama-3.1-8B & az \arm{E2x} ellenőrzője~\cite{grattafiori2024llama3}; \code{llama3.1:8b} & \tbd{[várhatóan Llama 3.1 Community License]} \\
gpt-oss-120b & bíró~\cite{gptoss2025}; egyetemi GenAI-szerver & \tbd{[várhatóan Apache 2.0]} \\
szervermodell & \arm{B-doc-S}, \arm{E4}, \arm{E4b}; egyetemi GenAI-szerver (\code{qwen38udq8xl}) & \tbd{[a pontos modell és licence]} \\
multilingual-e5-base & beágyazás~\cite{wang2024me5}; \nolinkurl{intfloat/multilingual-e5-base} & \tbd{[várhatóan MIT]} \\
\bottomrule
\end{tabularx}
\end{table}

A szervermodell és a bíró az egyetem saját, nem nyilvános infrastruktúráján futott, így a dokumentumok nem
kerültek külső szolgáltatóhoz, a két modell pontos változata viszont kívülről nem reprodukálható; a bírói
válaszok a \code{judge\_cache.jsonl} fájlokban megmaradtak, ezért az újrapontozás és az elemzés a szerver
nélkül is megismételhető. A kód MIT-licenc alatt érhető el \tbd{[a nyilvános tároló címe]}.
